\begin{abstract}
System-level network simulators such as ns-3 model the protocol stack with high fidelity but rely on analytical propagation models -- Friis, log-distance, 3GPP TR~38.901 stochastic -- that cannot reproduce environment-specific effects critical to 5G/6G physical-layer studies. Conversely, modern differentiable ray tracers such as Sionna~RT capture site-specific multipath, polarization, and antenna geometry but lack a discrete-event protocol stack. We present \emph{sionnart}, an ns-3 contrib module that embeds Sionna~RT directly into the simulator's address space via pybind11 and exposes it through standard \texttt{PropagationLossModel}, \texttt{PropagationDelayModel}, and \texttt{SpectrumPropagationLossModel} interfaces. A coherence-aware cache validated by a wavelength-and-array-derived displacement bound $\delta = \max(\delta_\mathrm{min}, \alpha\, d \cdot 0.886 / N_\mathrm{cols})$, coupled with a proactive virtual receiver generation mechanism that samples future channel states along mobile trajectories, suppresses redundant ray-tracer invocations while preserving channel realism. We evaluate sionnart in two campaigns: (i)~against the prior socket-coupled ns3sionna and a pure-ns-3 Friis baseline on an 802.11ax scenario across station counts $N \in \{1,\ldots,32\}$ and three mobility/load regimes; sionnart is up to ${\sim}232\times$ faster than ns3sionna at $N{=}32$ stationary load and ${\sim}44\times$ faster under high mobility while maintaining a minimal, flat GPU footprint under 500~MiB. It is even ${\sim}40\times$ faster than pure ns-3 in the stationary regime because cache hits replace recomputation. (ii)~On a 5G LENA NR setup with planar arrays $\{2{\times}2, 4{\times}4, 8{\times}8\}$ across free-space, urban-micro, and urban-macro scenes under three traffic loads. We demonstrate that deterministic multipath fading enforces more conservative, realistic CQI/MCS scheduling than stochastic models. This channel realism significantly shifts downstream metrics, altering gNB power consumption by up to $-29\%$ and uplink throughput by up to $+55\%$ relative to the analytical baseline, highlighting critical throughput and energy consumption trade-offs when scaling array sizes in complex environments. The artifact is a reusable contrib module and a methodology for fast, repeatable ray-traced network experiments suitable for digital-twin applications.
\end{abstract}

\begin{IEEEkeywords}
ns-3, Sionna~RT, ray tracing, 5G NR, energy modeling, propagation cache, network simulation, digital twin.
\end{IEEEkeywords}
